\section{Related Work and Discussion}

\noindent{\textbf{Attacks.}}
Several works have focused on systematically characterizing prompt injection attacks.
For instance, \cite{greshake2023not} demonstrate how hidden instructions can be inserted into external content to trigger harmful model behaviors.
Other works~\cite{wei2023jailbroken, zou2023universaltransferableadversarialattacks, zhang2024effective, deng2023masterkey, huang2024catastrophic, liu2024autodan, shen2024doanything, deng2024multilingual, chao2025jailbreaking} further contribute to our understanding by exploring a variety of prompt injection scenarios.
In addition,~\cite{Yang2024SOS} propose \texttt{SOS}, a training-time \say{soft prompt} attack that implants backdoors into open-source LLMs. The LLM behaves normally, until a trigger token activates malicious behaviors like jailbreaks, prompt stealing, or output manipulation.
In contrast,~\cite{russinovich2025greatwritearticlethat} introduce \texttt{Crescendo}, an inference-time jailbreak attack that unfolds over multiple dialogue turns.
Supporting this observation,~\cite{bullwinkel2025representation} performs an in-depth analysis of such multi-turn jailbreak attacks and show why defenses designed for single-turn interactions, like circuit breakers \cite{zou2024improving}, may be insufficient.
By examining internal model representations,~\cite{bullwinkel2025representation} shows that multi-turn prompts can gradually reframe harmful outputs as safe, effectively evading safeguards designed for single-turn interactions.
\texttt{BIPIA}~\cite{yi2025benchmarking} systematically assesses the risks posed by malicious instructions embedded in third-party content and introduces two defense paradigms: one based on prompt engineering and in-context strategies in a black-box setting, and another using adversarial training with special tokens in a white-box setting.
Finally,~\cite{abdelnabi2025llmail} introduces a realistic testbed for indirect prompt injection attacks targeting an LLM-based email assistant.
Collectively, these studies~\cite{greshake2023not, wei2023jailbroken, zou2023universaltransferableadversarialattacks, zhang2024effective, deng2023masterkey, huang2024catastrophic, liu2024autodan, shen2024doanything, deng2024multilingual, chao2025jailbreaking, Yang2024SOS, russinovich2025greatwritearticlethat, bullwinkel2025representation, yi2025benchmarking, abdelnabi2025llmail} provide a comprehensive overview of the diverse strategies employed to exploit prompt injection attacks, underscoring the ongoing need for robust security measures in LLM-integrated applications and services.


\noindent{\textbf{Attack Defenses.}}
Recent work proposes a variety of defense strategies against prompt injection attacks. These approaches differ in their methods and operating settings, ranging from inference-time countermeasures to white-box techniques and architectural safeguards.
\texttt{RigorLLM}~\cite{yuan2024rigorllm} framework enhances content moderation. Hines et al.~\cite{Hines2024Defending} build on the work of ~\cite{yi2025benchmarking}
to address indirect prompt injection attacks by proposing a defense strategy based on \emph{spotlighting}.
\texttt{CaMeL}~\cite{Debenedetti2025defeating} is a defense mechanism that explicitly detects and prevents prompt injection attacks by designing a secure execution layer around LLM operations.
The work sits at the intersection of generative AI and software security, bringing established principles from the latter to protect emerging systems.
SafeEar~\cite{SafeEar} proposes a method for detecting deepfake audio without accessing the actual speech content.
\texttt{TaskTracker}~\cite{abdelnabi2025get} is a white-box defense,  that centers on tracking the model’s internal activations before and after processing external text. 
\texttt{DataSentinel}~\cite{liu2025datasentinel} detects prompt injection attacks by fine-tuning a language model with a game-theoretic mini-max optimization. It uses a detection instruction embedding a secret key to verify the integrity of incoming data.
Beurer-Kellner et al.~\cite{beurerkellner2025design} propose a six architectural design patterns (action-selector, plan-then-execute, map-reduce, dual LLM, code-then-execute, and context-minimization) that secure LLM agents against attacks. 
Overall, most existing defenses are probabilistic and can be bypassed~\cite{Debenedetti2025defeating, costa2025securingaiagentsinformationflow}. Retrospective systems are needed to identify attackers, yet cross-service threat intelligence for large-scale prompt injection attacks correlation remains unexplored. To address this gap, techniques like SimHash can be leveraged. However, SimHash captures only syntactic similarity and fails to capture semantic similarity (Figures~\ref{fig:binaryshield_results},~\ref{fig:simhash_results}, and~\ref{fig:simhash_vs_binaryshield_f1}). 
\emph{\tool is the first to focus on this gap, proposing a concrete system for cross-service threat intelligence and paving the way for future research in this domain.}








\subsection{Discussion}  
\label{sec:discussion-and-threats-to-validity}

Our study has limitations similar to those in any research. 
Synthetic paraphrasing and word substitution may yield varying results with different LLMs and prompts. 
The effectiveness of the \tool's PII redaction module depends on the accuracy of detection methods. 
Residual PII may remain in prompts.
%
Our results rely on specific embedding models, ModernBert (open-source)~\cite{warner-etal-2025-smarter, gill2025advancingsemanticcachingllms} and OpenAI text-embedding-3-large (proprietary). 
Thus, performance may vary with other models. 
We determined optimal Hamming distance thresholds through search on our datasets. 
However, these thresholds may not generalize to other attack distributions or organizational contexts.
%
All experiments contain English prompts. 
Future work should explore cross-lingual attack detection and domain-specific technical prompts. 
We depend on computational non-invertibility and differential privacy guarantees. 
The evaluation does not assess resistance to reconstruction attacks, which represent an orthogonal concern to current focus of the paper.


\noindent{\textbf{Differential Privacy.}} Noise can be added to the fingerprint through various differential privacy mechanisms~\cite{warner1965randomized, Wang2017Locally, bhowmick2018protection, rappor, feyisetan2021private, vadhan2021concurrent, carvalho2021brr, bollegala2023neighbourhood, xie2024differentially, meehan2022sentence, asi2022optimal, du2023sanitizing, bittau2017prochlo, fanti2016building, hsu2014differential, near_abuah_2021, fioretto2024differential, dwork2014algorithmic, pmlr-v80-balle18a, NIPS2017_253614bb, hu2024sok}. 
%
\tool utilizes foundational work of Randomized Response mechanism~\cite{warner1965randomized}. 
Assessing the impact of other DP mechanisms and global differential privacy on the overall privacy guarantees is an important direction for future research. 
We did not formulate the privacy utility trade off as an optimization problem. 
Readers may consult existing literature on this topic~\cite{near_abuah_2021, Wang2017Locally, dwork2014algorithmic, fioretto2024differential}. 
However, in our experiments, we empirically evaluated conditions ranging from extreme privacy, where accuracy approaches zero, to moderate and weak privacy.
Each organization has specific privacy policies, and we do not recommend a particular differential privacy budget. 
Organizations should work with their privacy teams to select an appropriate budget based on their guidelines.  
